\chapter{Upper Abdominal Pain}

\section*{Definition}
Upper abdominal pain occurs between the ribs and the navel and may arise from the stomach, liver, gallbladder, pancreas, bowel, heart, lungs, or abdominal wall.

\section*{When to See a Doctor}
Seek urgent evaluation for Severe sudden pain, rigid abdomen, vomiting blood, black stool, jaundice, fever, fainting, chest symptoms, or persistent vomiting. Arrange routine or prompt assessment for persistent, recurrent, unexplained, or function-limiting symptoms.

\section*{How It Develops}
Upper-abdominal pain may arise from the stomach, gallbladder, liver, pancreas, bowel, abdominal wall, or referred chest disease. Location, relation to meals, severity, vomiting, jaundice, fever, bleeding, and cardiopulmonary symptoms guide urgency and testing.

\section*{Causes}
\begin{itemize}
  \item Peptic disease, gastritis, gallbladder disease, pancreatitis, hepatitis, bowel disease, pneumonia, myocardial ischemia, and muscle strain.
  \item Medication effects, environmental exposures, and chronic disease may contribute.
  \item Less common but important causes include malignancy, vascular disease, or organ failure when symptoms persist or progress.
\end{itemize}

\section*{Related Symptoms}
\begin{itemize}
  \item Pain, fever, fatigue, nausea, or changes in normal organ function
  \item Bleeding, weight change, shortness of breath, weakness, or neurologic symptoms when a serious cause is present
\end{itemize}

\section*{Medical Evaluation}
\begin{itemize}
  \item History addresses onset, duration, severity, triggers, exposures, medications, medical conditions, pregnancy status when relevant, and warning symptoms.
  \item Examination focuses on vital signs and the organ systems suggested by the symptom.
  \item Testing may include blood or urine studies, cultures, imaging, physiologic testing, or specialist examination according to the clinical pattern.
\end{itemize}

\section*{Treatment Options}
Treatment depends on the diagnosis and may include supportive care, acid-related treatment, management of gallbladder or pancreatic disease, or urgent surgical or medical care. Avoid masking severe or progressive pain with repeated self-treatment.

\section*{Self-Care and Home Management}
Avoid known triggers, maintain reasonable hydration, record timing and associated symptoms, and use only treatments appropriate for the suspected cause. Do not use leftover antibiotics or delay urgent care because symptoms temporarily improve.

\section*{References}
\begin{thebibliography}{99}
\bibitem{upper_abdominal_pain:dyspepsia} Moayyedi P, Lacy BE, Andrews CN, et al. ACG and CAG Clinical Guideline: Management of Dyspepsia. \textit{Am J Gastroenterol}. 2017;112:988--1013.
\bibitem{upper_abdominal_pain:tokyo} Yokoe M, Hata J, Takada T, et al. Tokyo Guidelines 2018: diagnostic criteria and severity grading of acute cholecystitis. \textit{J Hepatobiliary Pancreat Sci}. 2018;25:41--54.
\bibitem{upper_abdominal_pain:harrison} Jameson JL, et al. \textit{Harrison's Principles of Internal Medicine}. 21st ed. McGraw-Hill; 2022.
\end{thebibliography}
